\documentclass[a4paper,10pt,BCOR10mm,oneside,headsepline]{scrartcl}
\usepackage[portuges]{babel}
\usepackage[utf8]{inputenc}
\usepackage{amsmath}
\usepackage{wasysym}% provides \ocircle and \Box
\usepackage{enumitem}% easy control of topsep and leftmargin for lists
\usepackage{color}% used for background color
\usepackage{forloop}% used for \Qrating and \Qlines
\usepackage{ifthen}% used for \Qitem and \QItem
\usepackage{typearea}
\usepackage{graphicx}
\usepackage{booktabs}
\usepackage{listings}
\usepackage{color}
\usepackage{algorithm}
\usepackage{algorithmic}
\usepackage{amssymb}
\usepackage{minted}
\usepackage{mdframed}
\usepackage{wasysym}
\usepackage{xcolor}

\usepackage{tgtermes}

\areaset{17cm}{26cm}
\setlength{\topmargin}{-2cm}

\definecolor{codegreen}{rgb}{0,0.6,0}
\definecolor{codegray}{rgb}{0.5,0.5,0.5}
\definecolor{codepurple}{rgb}{0.58,0,0.82}
\definecolor{backcolour}{rgb}{0.95,0.95,0.92}
 
\lstdefinestyle{mystyle}{
    backgroundcolor=\color{backcolour},   
    commentstyle=\color{codegreen},
    keywordstyle=\color{magenta},
    numberstyle=\tiny\color{codegray},
    stringstyle=\color{codepurple},
    basicstyle=\footnotesize,
    breakatwhitespace=false,         
    breaklines=true,                 
    captionpos=b,                    
    keepspaces=true,                 
    numbers=left,                    
    numbersep=5pt,                  
    showspaces=false,                
    showstringspaces=false,
    showtabs=false,                  
    tabsize=2
}
 
\lstset{style=mystyle}

\mdfsetup{
    skipabove=12pt,
    skipbelow=12pt,
    innertopmargin=8pt,
    innerbottommargin=8pt,
    innerleftmargin=10pt,
    innerrightmargin=10pt,
    splittopskip=12pt,
    splitbottomskip=12pt,
}

\algsetup{linenodelimiter=.,linenosize=\tiny}

\renewcommand{\algorithmicrequire}{\textbf{Entrada:}}
\renewcommand{\algorithmicensure}{\textbf{Saída:}}
\floatname{algorithm}{Algoritmo}

%% \Qq = Questionaire question. Oh, this is just too simple. It helps
%% making it easy to globally change the appearance of questions.
\newcommand{\Qq}[1]{\textbf{#1}}
%\newcommand{\Qq}[1]{{#1}}

%% \QO = Circle or box to be ticked. Used both by direct call and by
%% \Qrating and \Qlist.
\newcommand{\QO}{$\Box$}% or: $\ocircle$

%% \Qrating = Automatically create a rating scale with NUM steps, like
%% this: 0--0--0--0--0.
\newcounter{qr}
\newcommand{\Qrating}[1]{\QO\forloop{qr}{1}{\value{qr} < #1}{---\QO}}

%% \Qline = Again, this is very simple. It helps setting the line
%% thickness globally. Used both by direct call and by \Qlines.
\newcommand{\Qline}[1]{\noindent\rule{#1}{0.6pt}}

%% \Qlines = Insert NUM lines with width=\linewith. You can change the
%% \vskip value to adjust the spacing.
\newcounter{ql}
\newcommand{\Qlines}[1]{\forloop{ql}{0}{\value{ql}<#1}{\vskip0em\Qline{\linewidth}}}

%% \Qlist = This is an environment very similar to itemize but with
%% \QO in front of each list item. Useful for classical multiple
%% choice. Change leftmargin and topsep accourding to your taste.
\newenvironment{Qlist}{%
\renewcommand{\labelitemi}{\QO}
\begin{itemize}[leftmargin=1.5em,topsep=-.5em]
}{%
\end{itemize}
}

%% \Qtab = A "tabulator simulation". The first argument is the
%% distance from the left margin. The second argument is content which
%% is indented within the current row.
\newlength{\qt}
\newcommand{\Qtab}[2]{
\setlength{\qt}{\linewidth}
\addtolength{\qt}{-#1}
\hfill\parbox[t]{\qt}{\raggedright #2}
}

%% \Qitem = Item with automatic numbering. The first optional argument
%% can be used to create sub-items like 2a, 2b, 2c, ... The item
%% number is increased if the first argument is omitted or equals 'a'.
%% You will have to adjust this if you prefer a different numbering
%% scheme. Adjust topsep and leftmargin as needed.
\newcounter{itemnummer}
\newcommand{\Qitem}[2][]{% #1 optional, #2 notwendig
\ifthenelse{\equal{#1}{}}{\stepcounter{itemnummer}}{}
\ifthenelse{\equal{#1}{a}}{\stepcounter{itemnummer}}{}
\begin{enumerate}[topsep=2pt,leftmargin=2.8em]
\item[\textbf{\arabic{itemnummer}#1.}] #2
\end{enumerate}
}

%% \QItem = Like \Qitem but with alternating background color. This
%% might be error prone as I hard-coded some lengths (-5.25pt and
%% -3pt)! I do not yet understand why I need them.
\definecolor{bgodd}{rgb}{0.8,0.8,0.8}
\definecolor{bgeven}{rgb}{0.9,0.9,0.9}
\newcounter{itemoddeven}
\newlength{\gb}
\newcommand{\QItem}[2][]{% #1 optional, #2 notwendig
\setlength{\gb}{\linewidth}
\addtolength{\gb}{-5.25pt}
\ifthenelse{\equal{\value{itemoddeven}}{0}}{%
\noindent\colorbox{bgeven}{\hskip-3pt\begin{minipage}{\gb}\Qitem[#1]{#2}\end{minipage}}%
\stepcounter{itemoddeven}%
}{%
\noindent\colorbox{bgodd}{\hskip-3pt\begin{minipage}{\gb}\Qitem[#1]{#2}\end{minipage}}%
\setcounter{itemoddeven}{0}%
}
}

\newcommand{\rawcode}[1]{\texttt{#1}}
\newcommand{\code}[2]{\mintinline{#1}{#2}}
\newcommand{\file}[1]{\texttt{\underline{#1}}}

\renewcommand{\theFancyVerbLine}{\tiny \texttt{\textcolor{gray}{\arabic{FancyVerbLine}}}}

\newcommand{\inputcode}[2]{\begin{mdframed}\inputminted[fontsize=\footnotesize,linenos,numbersep=4pt]{#1}{#2}\end{mdframed}}

\newenvironment{solution}{%
    \begin{mdframed}%
        \noindent\textbf{Solu\c{c}\~ao:}\par\medskip
}{%
    \end{mdframed}%
}

\newcommand{\inputsnippet}[4]{\begin{mdframed}[shadow=true,shadowsize=3pt,shadowcolor=gray]\inputminted[fontsize=\footnotesize,linenos,firstline=#2,lastline=#3,numbersep=2pt]{#1}{#4}\end{mdframed}}

\newcommand{\inputsyntax}[2]{\inputminted[fontsize=\footnotesize]{#1}{#2}}


\begin{document}

\begin{center}
	\includegraphics[scale=0.45]{fga}
	% fga.ps: 1087x103 pixel, 72dpi, 38.35x3.63 cm, bb=0 0 1087 103
\end{center}


\vspace{0.1in}

\begin{center}
    \textbf{\huge Paradigmas de Programação}

    \vspace{0.1in}
    \textbf{\Large Cálculo-$\lambda$: Lista 1}
\end{center}\vskip1em

\renewcommand{\Qq}[1]{{#1}}

\section*{Parte A - Definição de cálculo-$\lambda$}

\Qitem{ \Qq{Descreva a sequência de aplicações dos axiomas de construção que forma o termo-$\lambda$
indicado. Assuma que $x, y, z, w$ são variáveis e siga o modelo do exemplo, o qual ilustra a construção do termo $wz(\lambda y.xy)$:
\[
    x 
    \quad \xrightarrow{\tiny \mbox{aplicação}}\quad xy
    \quad \xrightarrow{\tiny \mbox{abstração}}\quad \lambda y.xy
    \quad \xrightarrow{\tiny \mbox{aplicação}}\quad z(\lambda y.xy)
    \quad \xrightarrow{\tiny \mbox{aplicação}}\quad wz(\lambda y.xy)
\]
    \begin{Qlist}
    \item[(a)] $(\lambda y.(\lambda x.xy)xy)$
    \item[(b)] $\lambda w.xyx(\lambda z.z)$
    \end{Qlist}
} }

\begin{solution}
\textbf{(a)}
\[
    x
    \;\xrightarrow{\tiny\mbox{aplicação}}\; xy
    \;\xrightarrow{\tiny\mbox{abstração}}\; \lambda x.xy
    \;\xrightarrow{\tiny\mbox{aplicação}}\; (\lambda x.xy)x
    \;\xrightarrow{\tiny\mbox{aplicação}}\; (\lambda x.xy)xy
    \;\xrightarrow{\tiny\mbox{abstração}}\; \lambda y.(\lambda x.xy)xy
\]

\textbf{(b)} O termo $\lambda w.xyx(\lambda z.z)$ requer a construção de dois
subtermos, $xyx$ e $\lambda z.z$, que são então combinados por aplicação
e juntos por abstração em $w$:
\begin{align*}
    &z \;\xrightarrow{\tiny\mbox{abstração}}\; \lambda z.z \\[4pt]
    &x \;\xrightarrow{\tiny\mbox{aplicação}}\; xy
      \;\xrightarrow{\tiny\mbox{aplicação}}\; xyx
      \;\xrightarrow{\tiny\mbox{aplicação}}\; xyx(\lambda z.z)
      \;\xrightarrow{\tiny\mbox{abstração}}\; \lambda w.xyx(\lambda z.z)
\end{align*}
\end{solution}

\Qitem{ \Qq{Para cada termo abaixo, indique se a substituição da variável $x$ por $N$ se aplica ou
        não. Em caso afirmativo, realize a substituição; caso contrário, justifique porque a 
        substituição não se aplica.

    \begin{Qlist}
    \item[(a)] $(\lambda y.xyx)(yxy)$
    \item[(b)] $\lambda x.(\lambda y.xyx)$
    \item[(c)] $\lambda z.yzwz$
    \end{Qlist}
} }

\begin{solution}
\begin{Qlist}
\item[(a)] \textbf{Aplica-se.} A variável $x$ ocorre livre em ambos os subtermos.
Aplicando as regras de substituição para aplicação e abstração
(supondo $y\notin FV(N)$ para evitar captura no contexto de $\lambda y$):
\[
    \bigl((\lambda y.xyx)(yxy)\bigr)[x:=N]
    = (\lambda y.xyx)[x:=N]\cdot(yxy)[x:=N]
    = (\lambda y.NyN)(yNy).
\]

\item[(b)] \textbf{Não se aplica.} O símbolo $\lambda x$ mais externo ata a variável
$x$ em todo o corpo do termo; portanto, não há ocorrência \emph{livre} de $x$ sobre
a qual a substituição possa agir.

\item[(c)] \textbf{Aplica-se, mas não altera o termo.} A variável $x$ não ocorre
no termo $\lambda z.yzwz$; pelas regras ii e iv (como no exemplo
$(\lambda x.xy)[z:=N]=\lambda x.xy$ dos slides):
\[
    (\lambda z.yzwz)[x:=N] = \lambda z.\,(yzwz)[x:=N] = \lambda z.yzwz.
\]
\end{Qlist}
\end{solution}

\Qitem{ \Qq{Aplique os axiomas de redução, enquanto for possível, nas expressões abaixo:
    \begin{Qlist}
    \item[(a)] $({\lambda y.({\lambda x.xyx})})x$
    \item[(b)] $\lambda z.({\lambda y.yz})x$
    \end{Qlist}
} }

\begin{solution}
\begin{Qlist}
\item[(a)] Aplicamos a redução-$\beta$ substituindo $y$ por $x$ no corpo $(\lambda x.xyx)$:
\[
    (\lambda y.(\lambda x.xyx))\,x \;\to_\beta\; (\lambda x.xyx)[y:=x].
\]
O termo substituto $x$ tem $x\in FV(x)$ e $y$ ocorre livre no corpo $xyx$ de $\lambda x$;
logo a substituição direta faria $\lambda x$ capturar o $x$ livre do argumento. Usamos a conversão-$\alpha$ para renomear a variável atada:
\[
    \lambda x.xyx \;\equiv_\alpha\; \lambda z.zyz.
\]
Como $z\neq x$, a substituição não causa captura:
\[
    (\lambda z.zyz)[y:=x] = \lambda z.\,(zyz)[y:=x] = \lambda z.zxz.
\]
O resultado $\lambda z.zxz$ é uma forma normal; nenhuma redução adicional é possível.

\item[(b)] Aplicamos a redução-$\beta$ à subexpressão interna $(\lambda y.yz)x$:
\[
    \lambda z.\,(\lambda y.yz)\,x
    \;\to_\beta\;
    \lambda z.\,(yz)[y:=x]
    = \lambda z.xz.
\]
O termo $\lambda z.xz$ tem a forma $\lambda z.Mz$ com $M=x$ e
$z\notin FV(x)=\{x\}$; a conversão-$\eta$ é, portanto, aplicável:
\[
    \lambda z.xz \;\equiv_\eta\; x.
\]
O resultado final é $x$.
\end{Qlist}
\end{solution}

\Qitem{ \Qq{Identifique, para cada termo abaixo, se a conversão-$\eta$ é aplicável. Justifique caso
        não seja, ou aplique a conversão, se possível.
        \begin{Qlist}
        \item[(a)] $\lambda x.x$
        \item[(b)] $\lambda x.y$
        \item[(c)] $\lambda x.xyx$
        \end{Qlist}
} }

\begin{solution}
A conversão-$\eta$ transforma $\lambda x.Mx$ em $M$. Ela só se aplica quando o
corpo da abstração é uma \emph{aplicação} da forma $M\,x$ com $x\notin FV(M)$.

\begin{Qlist}
\item[(a)] $\lambda x.x$: O corpo é a variável $x$, e não uma aplicação da forma
$M\,x$. A conversão-$\eta$ \textbf{não é aplicável}.

\item[(b)] $\lambda x.y$: O corpo é a variável $y$; tampouco é uma aplicação.
A conversão-$\eta$ \textbf{não é aplicável}.

\item[(c)] $\lambda x.xyx$: O corpo $xyx=(xy)x$ está na forma $M\,x$ com $M=xy$.
No entanto, $x\in FV(xy)=\{x,y\}$, o que viola a condição $x\notin FV(M)$.
A conversão-$\eta$ \textbf{não é aplicável}.
\end{Qlist}
\end{solution}

\Qitem{ \Qq{Determine o conjunto das variáveis livres do termo $({\lambda x.zxyw})({\lambda y.wxzy})$.} }

\begin{solution}
Pelas regras de variáveis livres $FV$:
\begin{align*}
    FV(\lambda x.zxyw) &= FV(zxyw) - \{x\} = \{z,x,y,w\} - \{x\} = \{z,y,w\},\\
    FV(\lambda y.wxzy) &= FV(wxzy) - \{y\} = \{w,x,z,y\} - \{y\} = \{w,x,z\}.
\end{align*}
Logo:
\[
    FV\!\bigl((\lambda x.zxyw)(\lambda y.wxzy)\bigr)
    = \{z,y,w\}\cup\{w,x,z\} = \{x,\,y,\,z,\,w\}.
\]
\end{solution}

\Qitem{ \Qq{Escreva um combinador formado apenas por $S$ e $K$ que seja extensionalmente
        igual ao combinador $K_*$.} }

\begin{solution}
O combinador $K_*\equiv\lambda xy.y$ seleciona sempre o \emph{segundo} argumento,
e $SK$ se comporta da mesma forma:
\[
    SK\,x\,y
    \;=\; Ky\,(xy)
    \;=\; y
    \;=\; K_*\,x\,y,
\]
pois $Sxyz=xz(yz)$ e $Kab=a$. Portanto, $K_* = SK$.
\end{solution}

\section*{Parte B - Lógica booleana}

\Qitem{ \Qq{Defina um termo lambda $X$ que correponda à operação lógica ``ou exclusivo'', isto é,
        $Xpq = T$ se $p$ e $q$ tem valores lógicos distintos e $Xpq = F$, caso contrário.} }

\begin{solution}
Usando $T\equiv\lambda xy.x$, $F\equiv\lambda xy.y$ e
$\lnot q\equiv qFT$, definimos:
\[
    X \;\equiv\; \lambda pq.\;p\,(\lnot q)\,q
      \;=\; \lambda pq.\;p\,(qFT)\,q.
\]

\textbf{Verificação} (tabela-verdade):
\begin{center}
\begin{tabular}{cc|l|c}
\toprule
$p$ & $q$ & $X\,p\,q = p\,(qFT)\,q$ & esperado \\
\midrule
$T$ & $T$ & $T\underbrace{(TFT)}_{=F}T = TFT = F$ & $F$ \checkmark \\[4pt]
$T$ & $F$ & $T\underbrace{(FFT)}_{=T}F = TTF = T$ & $T$ \checkmark \\[4pt]
$F$ & $T$ & $F\underbrace{(TFT)}_{=F}T = FFT = T$ & $T$ \checkmark \\[4pt]
$F$ & $F$ & $F\underbrace{(FFT)}_{=T}F = FTF = F$ & $F$ \checkmark \\
\bottomrule
\end{tabular}
\end{center}
\end{solution}

\Qitem{ \Qq{Escreva, para cada operador lógico abaixo, um combinador na base $SK$ que seja
        extensionalmente igual ao operador. Indique a posição do operador (pré-fixada, in-fixada ou
        pós-fixada).
    \begin{Qlist}
    \item[(a)] condicional
    \item[(b)] bicondicional
    \end{Qlist}
} }

\begin{solution}
Com $T=K$ e $F=SK$:

\begin{Qlist}
\item[(a)] \textbf{Condicional $p\to q$.} Combinador: $K$, posição pós-fixada.

$p\,q\,K = p\to q$:
\begin{itemize}
    \item Se $p=T$: $\;T\,q\,K = q$ \quad (T seleciona o primeiro argumento).
    \item Se $p=F$: $\;F\,q\,K = K = T$ \quad (F seleciona o segundo argumento).
\end{itemize}

\textbf{Verificação} (tabela-verdade):
\begin{center}
\begin{tabular}{cc|c|c}
\toprule
$p$ & $q$ & $p\,q\,K$ & $p\to q$ \\
\midrule
$T$ & $T$ & $T\,T\,K = T$     & $T$ \checkmark \\
$T$ & $F$ & $T\,F\,K = F$     & $F$ \checkmark \\
$F$ & $T$ & $F\,T\,K = K = T$ & $T$ \checkmark \\
$F$ & $F$ & $F\,F\,K = K = T$ & $T$ \checkmark \\
\bottomrule
\end{tabular}
\end{center}
Portanto, $\text{COND} = K$ (pós-fixado): $p\,q\,K\;=\;p\to q$.

\item[(b)] \textbf{Bicondicional $p\leftrightarrow q$.} Combinador: $B$, posição pré-fixada.

$p\leftrightarrow q$ é verdadeiro quando $p$ e $q$ têm o mesmo valor: para $p=T$
o resultado é $q$, e para $p=F$ o resultado é $\lnot q$. O combinador pré-fixado
\[
    B \;\equiv\; \lambda p.\;p\,I\,\text{NOT},
\]
onde $I=SKK$ e $\text{NOT}\equiv\lambda q.\;q\,F\,T = \lambda q.\;q\,(SK)\,K$ (a versão
\emph{pré-fixada} da negação; o $\mathbf{NOT}\equiv(\mathbf{SK})\mathbf{K}$ dos slides é
\emph{pós-fixado}, isto é, $q\,\mathbf{NOT} = q\,(SK)\,K$, e não serviria aqui, pois
$(SK)K\,q = q$), satisfaz
\[
    B\,p\,q = (p\,I\,\text{NOT})\,q =
    \begin{cases} I\,q = q, & \text{se } p=T,\\ \text{NOT}\,q = \lnot q, & \text{se } p=F. \end{cases}
\]

\textbf{Verificação}:
\begin{center}
\begin{tabular}{cc|c|c}
\toprule
$p$ & $q$ & $B\,p\,q$ & $p\leftrightarrow q$ \\
\midrule
$T$ & $T$ & $I\,T = T$          & $T$ \checkmark \\
$T$ & $F$ & $I\,F = F$          & $F$ \checkmark \\
$F$ & $T$ & $\lnot T = F$       & $F$ \checkmark \\
$F$ & $F$ & $\lnot F = T$       & $T$ \checkmark \\
\bottomrule
\end{tabular}
\end{center}

\textbf{Derivação em $SK$ puro}:
\[
    B = \lambda p.\;p\,I\,\text{NOT}
      = S(\lambda p.\;p\,I)(\lambda p.\;\text{NOT}).
\]
Como $\text{NOT}$ não depende de $p$:
$\lambda p.\,\text{NOT} = K\,\text{NOT}$. Traduzindo $\text{NOT}$ para $SK$:
\[
    \text{NOT} = \lambda q.\;q\,(SK)\,K
               = S(\lambda q.\;q\,(SK))(\lambda q.\;K)
               = S\bigl(S\,I\,(K(SK))\bigr)(KK)
               = S\bigl(S(SKK)(K(SK))\bigr)(KK).
\]
(Verificação: $S(SI(K(SK)))(KK)\,q = SI(K(SK))\,q\,(KKq) = Iq\,(K(SK)q)\,K = q\,(SK)\,K$.)

E, usando $I=SKK$:
\[
    \lambda p.\;p\,I = S(\lambda p.p)(\lambda p.I) = S\,I\,(K\,I)
                    = S(SKK)(K(SKK)).
\]
Portanto:
\[
    B = S\bigl(S(SKK)(K(SKK))\bigr)\Bigl(K\bigl(S(S(SKK)(K(SK)))(KK)\bigr)\Bigr).
\]
\end{Qlist}
\end{solution}

\end{document}
